\documentclass[11pt]{article}
\usepackage[margin=0.9in, top=0.75in, bottom=0.75in]{geometry}
\usepackage[utf8]{inputenc}
\usepackage[T1]{fontenc}
\usepackage{hyperref}
\usepackage{url}
\usepackage{booktabs}
\usepackage{amsmath}
\usepackage{graphicx}
\usepackage{subcaption}
\usepackage{enumitem}
\usepackage{parskip}
\usepackage{microtype}
\usepackage{titlesec}

\setlength{\parskip}{3pt}
\setlength{\parindent}{0pt}

\titlespacing*{\section}{0pt}{6pt}{2pt}
\titlespacing*{\subsection}{0pt}{4pt}{1pt}
\titleformat{\section}{\normalfont\bfseries\large}{}{0em}{}
\titleformat{\subsection}{\normalfont\bfseries}{}{0em}{}

\begin{document}

\begin{center}
  {\large\bfseries Replicating DreamBooth: Subject-Driven Fine-Tuning of Stable Diffusion}\\[4pt]
  Ryan Qiu \quad Gordon Mei \quad Caleb Shim \quad Evan Cui \\
  Cornell University \quad CS4782 --- Deep Learning
\end{center}

\vspace{2pt}\hrule\vspace{6pt}

%% ─── 1. INTRODUCTION ───────────────────────────────────────────────────────────
\section{Introduction}

Large-scale text-to-image diffusion models synthesize photorealistic images from prompts but cannot render a \emph{specific} subject unseen during training. \textbf{DreamBooth}~\cite{ruiz2023dreambooth} (Ruiz et al., CVPR 2023) addresses this by fine-tuning the model on 3--5 reference images, binding a particular subject to a rare identifier token and enabling compositional generation in novel contexts.

We replicate DreamBooth on Stable Diffusion v1.5 (SD~1.5)~\cite{rombach2022ldm}, implementing the full training pipeline from scratch. Our goal is to reproduce the paper's reported quantitative metrics on SD~1.5 and document the implementation details necessary to achieve them.

%% ─── 2. CHOSEN RESULT ──────────────────────────────────────────────────────────
\section{Chosen Result}

We target the quantitative evaluation from Table~2 of the DreamBooth paper, which reports subject fidelity (DINO~\cite{caron2021dino}, CLIP-I~\cite{radford2021clip}) and prompt fidelity (CLIP-T) for fine-tuned SD~1.5 models across 30 subjects. The paper reports mean scores of DINO~0.668, CLIP-I~0.803, and CLIP-T~0.305 for SD. This result is central to the paper's contribution: it demonstrates that subject identity and prompt-following can coexist in a fine-tuned diffusion model. Reproducing it validates that our pipeline correctly implements both the rare-token binding and prior preservation mechanisms.

%% ─── 3. METHODOLOGY ────────────────────────────────────────────────────────────
\section{Methodology}

\textbf{Architecture.} SD~1.5 has three components: a VAE (${\approx}$83M params) that encodes images to latents, a UNet (${\approx}$860M params) that denoises in latent space, and a CLIP text encoder (${\approx}$123M params). We fine-tune the UNet only; both the VAE and CLIP text encoder are frozen throughout training.

\textbf{Training objective.} Instance images are captioned as ``a photo of \texttt{sks} [class]'', where \texttt{sks} is a rare identifier with negligible prior semantics. At each step, both instance and class images are encoded to latents, noise is sampled and added, and the UNet predicts the noise. The combined loss is:
\[
  \mathcal{L} =
  \underbrace{\|\hat{\varepsilon}_\text{inst} - \varepsilon_\text{inst}\|^2}_{\text{instance loss}}
  + \lambda\,
  \underbrace{\|\hat{\varepsilon}_\text{class} - \varepsilon_\text{class}\|^2}_{\text{prior preservation}},
  \quad \lambda = 1.0.
\]
Prior-preservation class images (100 per subject) are generated by the unmodified SD~1.5 model before training, substituting for the Google Imagen model used in the original paper.

\textbf{Hyperparameters.} 800 steps, learning rate $5{\times}10^{-6}$, resolution $512{\times}512$, batch size~1, 8-bit AdamW optimizer~\cite{dettmers2022llm}, gradient checkpointing. The UNet runs in FP32 while frozen components use FP16.

\textbf{Dataset and evaluation.} We fine-tune on \texttt{dog}, \texttt{dog2}, and \texttt{backpack} from the official DreamBooth dataset~\cite{dreambooth_dataset}, each with 5 instance images. Evaluation generates 4 images per prompt across 5 prompts (20 images/subject). DINO uses ViT-S/16~\cite{caron2021dino}; CLIP-I and CLIP-T use CLIP ViT-L/14~\cite{radford2021clip}. Training runs on a Google Colab T4 GPU (${\approx}$35 min/subject); metrics are computed on CPU.

%% ─── 4. RESULTS & ANALYSIS ─────────────────────────────────────────────────────
\section{Results \& Analysis}

\begin{table}[h]
  \centering
  \caption{Quantitative metrics. ``Real images'' is the upper bound; ``Paper (SD)'' is the DreamBooth paper's SD~1.5 result.}
  \label{tab:metrics}
  \vspace{2pt}
  \small
  \begin{tabular}{lccc}
    \toprule
    & \textbf{DINO}$\uparrow$ & \textbf{CLIP-I}$\uparrow$ & \textbf{CLIP-T}$\uparrow$ \\
    \midrule
    dog              & 0.716 & 0.867 & 0.275 \\
    dog2             & 0.542 & 0.828 & 0.247 \\
    backpack         & 0.536 & 0.887 & 0.274 \\
    \midrule
    \textbf{Ours (mean)} & \textbf{0.598} & \textbf{0.861} & \textbf{0.265} \\
    \midrule
    Real images      & 0.774 & 0.885 & --- \\
    Paper (SD)       & 0.668 & 0.803 & 0.305 \\
    \bottomrule
  \end{tabular}
\end{table}

Our mean DINO (0.598) falls slightly below the paper's 0.668, which is expected given that we evaluate only 3 subjects rather than the full 30. CLIP-I (0.861) exceeds the paper's 0.803, indicating strong semantic alignment between generated and reference images. CLIP-T (0.265) is below the paper's 0.305, reflecting the inherent subject--prompt fidelity trade-off: our fine-tuned models preserve identity well but sacrifice some text adherence. Qualitatively, generated images maintain subject appearance across novel backgrounds and lighting; see Figure~\ref{fig:qual}.

The most notable gap is DINO variance across subjects: \texttt{dog} scores 0.716 (near the paper's mean) while \texttt{dog2} and \texttt{backpack} score around 0.54, suggesting that fine-grained identity binding depends on inter-image consistency in the instance set. The backpack scores high on CLIP-I (0.887) despite lower DINO, indicating that semantic similarity is preserved even where exact feature identity is weaker.

%% ─── 5. REFLECTIONS ────────────────────────────────────────────────────────────
\section{Reflections}

\textbf{Prior preservation is load-bearing, not optional.} When we ablated prior preservation, training did not degrade gracefully---the subject displaced the model's entire class representation, and general generation became incoherent. The prior loss is not a soft regularizer; it is necessary for stable training.

\textbf{FP16 training of the UNet causes a silent failure.} Using half-precision for the UNet caused gradient updates to underflow, effectively stopping weight updates. The loss curve appeared normal and no exceptions were raised, but outputs showed no identity binding. Keeping the UNet in FP32 (while frozen components use FP16) resolved the issue immediately. This failure is invisible without knowing to check for it.

\textbf{Freezing the text encoder improves results.} Jointly fine-tuning the UNet and CLIP text encoder led to overfitting: outputs lost compositional flexibility and failed to follow novel prompts. The text encoder already encodes language structure; fine-tuning it allows the model to ``cheat'' by redefining token semantics rather than adjusting the visual pathway. Keeping it frozen constrains learning to where it matters.

\textbf{Extensions and limitations.} We applied DreamBooth to a course logo graphic (\texttt{CS4782}) and to artistic style images. Logo fine-tuning faithfully reproduced the graphic in direct generation but failed at context placement (e.g., ``a mug with a \texttt{sks} logo''), with the model rendering the identifier as literal glyphs rather than the learned graphic. Style transfer showed partial aesthetic transfer but inconsistent fidelity, highlighting that DreamBooth's subject-binding objective is poorly matched to distributed visual properties like style.

\textbf{Future work.} Scaling to 30 subjects would enable a direct benchmark comparison. Ablations over step count and learning rate could map the under/overfitting boundary. Multi-subject binding---associating multiple identifiers in a single fine-tune---is an extension not covered by our replication.

%% ─── REFERENCES ──────────────────────────────────────────────────────────────
\begin{thebibliography}{9}

\bibitem{ruiz2023dreambooth}
N.~Ruiz, Y.~Li, V.~Jampani, Y.~Pritch, M.~Rubinstein, and K.~Aberman.
\newblock DreamBooth: Fine tuning text-to-image diffusion models for subject-driven generation.
\newblock In \textit{CVPR}, 2023.

\bibitem{rombach2022ldm}
R.~Rombach, A.~Blattmann, D.~Lorenz, P.~Esser, and B.~Ommer.
\newblock High-resolution image synthesis with latent diffusion models.
\newblock In \textit{CVPR}, 2022.

\bibitem{radford2021clip}
A.~Radford et al.
\newblock Learning transferable visual models from natural language supervision.
\newblock In \textit{ICML}, 2021.

\bibitem{caron2021dino}
M.~Caron et al.
\newblock Emerging properties in self-supervised vision transformers.
\newblock In \textit{ICCV}, 2021.

\bibitem{dettmers2022llm}
T.~Dettmers, M.~Lewis, Y.~Belkada, and L.~Zettlemoyer.
\newblock LLM.int8(): 8-bit matrix multiplication for transformers at scale.
\newblock In \textit{NeurIPS}, 2022.

\bibitem{diffusers}
HuggingFace.
\newblock Diffusers: State-of-the-art diffusion models.
\newblock \url{https://github.com/huggingface/diffusers}, 2022.

\bibitem{dreambooth_dataset}
Google.
\newblock DreamBooth dataset.
\newblock \url{https://github.com/google/dreambooth}, 2022.

\end{thebibliography}

%% ─── FIGURES (excluded from 2-page limit) ──────────────────────────────────────
\newpage
\begin{figure}[h]
  \centering
  \begin{subfigure}[b]{0.28\linewidth}
    \centering
    \includegraphics[width=\linewidth]{../poster/imgs/dog_ref.jpg}
    \caption{Reference subject}
  \end{subfigure}
  \hspace{2em}
  \begin{subfigure}[b]{0.28\linewidth}
    \centering
    \includegraphics[width=\linewidth]{../poster/imgs/cornellmerch2.png}
    \caption{``a photo of [v] dog with Cornell merch''}
  \end{subfigure}
  \caption{Reference image and a generated image using the fine-tuned model, demonstrating subject identity preserved in a novel context.}
  \label{fig:qual}
\end{figure}

\end{document}
